Soit le couple d'oxydoréduction diiode\ttslash{}ion iodure~:
$\ce{I2_{(aq)}\ttslash{}I^-_{(aq)}}$.
\begin{enumerate}
    \item Le diiode $\ce{I2_{(s)}}$ est légèrement soluble dans l'eau avec
          laquelle il donne une solution jaunâtre.
    
          Dans une solution aqueuse de diiode, on verse une solution de
          thiosulfate de sodium~:
          \[
              \ce{2 Na^+_{(aq)} + S2O^2-_{3(aq)}}
          \]

          Qu'observe-t-on~? Écrire l'équation de la réaction d'oxydoréduction
          qui s'effectue.
    \item L'iodure de potassium $\ce{KI_{(s)}}$ est très soluble dans l'eau avec
          laquelle il donne une solution incolore de formule~:
          \[
              \ce{K^+_{(aq)} + I^-_{(aq)}}
          \]

          Dans une solution aqueuse d'iodure de potassium, on verse une quantité
          modérée d'une solution de permanganate de potassium
          $\ce{K^+_{(aq)} + MnO^-_{4(aq)}}$ acidifiée par de l'acide
          sulfurique.
    
          Qu'observe-t-on~? Écrire l'équation de la réaction d'oxydoréduction
          qui s'effectue.
    \item On donne les couples d'oxydoréduction~:
          \[
              \ce{SO^2-_{4(aq)}}\ttslash{}\ce{SO2_{(aq)}}
              \quad\text{et}\quad
              \ce{MnO^-_{4(aq)}}\ttslash{}\ce{Mn^2+_{(aq)}}
          \]
          (en milieu acide).

          Quand on introduit un papier imbibé d'une solution acidifiée de
          permanganate de potassium dans un flacon contenant du dioxyde de
          soufre, celui-ci est décoloré. Interpréter en écrivant l'équation de
          cette réaction.
\end{enumerate}

